\chapter{Measure Theory}
Let \(\scrB\) be the set of all Bernoulli sequences. \(\scrB\) is uncountable.
\begin{proposition}
    If we delete a countable subset of \(\scrB\), we can index what is left by the points on the real interval \(I = \opcl{0}{1}\). That is, there exists an injective function \(f: I \to \scrB\).
\end{proposition}
\begin{proof}
    Each \(\omega \in I\) can be written as 
    \begin{align*}
        \omega &= \sum_{i = 1}^{\infty} \dfrac{a_i}{2^i} \qquad a_i = 0,1\\
        &= 0.a_1a_2\dots
    \end{align*}
    Since \(\omega\) may not have a unique binary representation, we will only consider non-terminating expansion for \(\omega\). Then, by mapping \(1\) to \(H\) and \(0\) to \(T\) we get an injective function from \(I\) to \(\scrB\). To show this, suppose \(\scrB_{\deg}\) is the set of all Bernoulli sequences that after a certain point degenerates to all tails.
    \begin{lemma}
        \(\scrB_{\deg}\) is countable.
    \end{lemma}
    \begin{prooflemma}
        Let \(\scrB_{\deg}^k\) be all degenerate Bernoulli sequences where we have only tails after \(k_{\cardinalTH}\) toss. Then, \(\scrB_{\deg}^k\) is finite and 
        \begin{equation*}
            \scrB_{\deg} = \bigcup_{k = 1}^{\infty} \scrB_{\deg}^k
        \end{equation*}
        is the countable union of finite set and hence \(\scrB_{\deg}\) is countable.
    \end{prooflemma}

    This concludes the proof.
\end{proof}

\begin{definition}[Borel Principle]
    Suppose \(E\) is a probabilistic event occuring in certain sequences. Let \(\scrB_{E}\) denote the subset of \(\scrB\) for which that event occurs. Let \(I_E\) be the corresponding subset of \(I\), then 
    \begin{equation*}
        \prob{E} = \func{\mu_L}{I_E}
    \end{equation*}
    where \(\mu_L\) is Lebesgue measure.
\end{definition}

\begin{example}
    Start with \(X\) dollars and at each toss you win \(1\) dollars if heads shows up and ypu lose \(1\) dollars if tail shows up. What is the probability you lose all your original stake?

    For \(\omega \in I\) define the \(k_{\cardinalTH}\) \textit{Radamcher} function, \(R_k\), by 
    \begin{align*}
        \func{R_k}{\omega} &= 2a_k - 1\\
        &= \begin{cases}
            +1 & a_k = 1\\
            -1 & a_k = 0
        \end{cases}
    \end{align*}
    Then, let \(\func{S_k}{\omega}\) be the total amount won or loss at \(k_{\cardinalTH}\) toss.
    \begin{equation*}
        \func{S_k}{\omega} = \sum_{l = 1}^k \func{R_l}{\omega}
    \end{equation*}
    Thus, the event that we lose our stake at \(k_{\cardinalTH}\) is 
    \begin{equation*}
        I_{E_k} = \set[[\Bigg]]<\omega \in I>{\func{S_l}{\omega} > -X \ \mathrm{for} \ l < k, \func{S_k}{\omega} = -X}
    \end{equation*}
    hence 
    \begin{equation*}
        I = \bigcup_{l = 1}^{\infty} I_{E_k}
    \end{equation*}
    is the event that we loss all our money eventually. We will however postpone calculating \(\func{\mu_L}{I_E}\) as it is not finite union of intervals.
\end{example}

\section{Weak law of large numbers}
For some fixed \(\epsilon\) let
\begin{equation*} 
    I_N = \set<\omega \in I>{\abs{\dfrac{\func{s_N}{\omega}}{N} - \dfrac{1}{2}} > \epsilon}
\end{equation*}
where \(\func{s_N}{\omega} = a_1 + a_2 + \dots + a_N\). Then, \(I_N\) represents the event that the number of heads and tails \underline{are not} roughly equal.
\begin{theorem}[Weak law of large numbers]
    WLLN states 
    \begin{equation*}
        \func{\mu_L}{I_N} \to 0 \ \mathrm{as} \ N \to \infty
    \end{equation*}
\end{theorem}
\begin{proof}
    Equivalently, for 
    \begin{equation*} 
        A_N = \set[[\Big]]<\omega \in I>{\abs{\func{S_N}{\omega}} > 2N \epsilon}
    \end{equation*}
    then
    \begin{equation*}
        \func{\mu_L}{A_N} \to 0  \ \mathrm{as} \ N \to \infty
    \end{equation*}
    \begin{lemma}[Chebyshev's inequality]
        Let \(f\) be a non-negative, piecewise constant function on \(\opcl{0}{1}\). Let \(\alpha > 0\) be given. Then, 
        \begin{equation*}
            \func{\mu_L}{\set[[\Big]]<\omega \in I>{\func{f}{\omega} > \alpha}} < \dfrac{1}{\alpha} \int_{0}^1 f \diffOperator x
        \end{equation*}
        where \(\int\) is the Riemann integral.
    \end{lemma}
    \begin{prooflemma}
        Since \(f\) is piecewise constant then there is a parition \(0 = x_1 < \dots < x_k = 1\) such that \(f = c_i\) on \(\opcl{x_i}{X_{i+1}}\) for \(i= 1, \dots k - 1\). Then, 
        \begin{align*}
            \int_{0}^1 f \diffOperator x &= \sum_{i = 1}^{k-1} c_i \bracket{x_{i+1} - x_i} \\
            &\geq  \sum_{c_i > \alpha} c_i \bracket{x_{i+1} - x_i} \\
            &> \alpha  \sum_{c_i > \alpha}  x_{i+1} - x_i = \alpha \func{\mu_L}{\set[[\Big]]<\omega \in I>{\func{f}{\omega} > \alpha}}
        \end{align*}
        which proves the Chebyshev's inequality.
    \end{prooflemma}
    Then, we have 
    \begin{equation*}
        A_N = \set[[\Big]]<\omega \in I>{\abs{\func{S_N}{\omega}} > 2N \epsilon}= \set[[\Big]]<\omega \in I>{\bracket{\func{S_N}{\omega}}^2 > 4N^2 \epsilon^2}
    \end{equation*}
    and hence 
    \begin{align*}
        \func{\mu_L}{A_N} &< \dfrac{1}{4N^2\epsilon^2} \int_0^1 \bracket{\func{S_N}{\omega}}^2 \diffOperator \omega\\
        &= \dfrac{1}{4N^2\epsilon^2} \squareBracket{\sum_i \int_0^1 \bracket{\func{R_i}{\omega}}^2 \diffOperator \omega + \sum_{i\neq j} \int_0^1 \func{R_i}{\omega}\func{R_j}{\omega} \diffOperator \omega }\\
        &= \dfrac{1}{4N^2\epsilon^2} N = \dfrac{1}{4N \epsilon^2}
    \end{align*}
    Therefore, as \(N\) approaches infinity, \(\func{\mu}{A_N}\) approaches zero.
\end{proof}

Now we want to show that for a ``typical'' Bernoulli sequence 
\begin{equation}\label{eq:typicalBernoulli}
    \dfrac{1}{2} - \dfrac{\func{s_N}{\omega}}{N} \to 0 \ \mathrm{as} \ N \to \infty
\end{equation}
and by ``typical'' we mean \ ref{eq:typicalBernoulli} fails on a set of zero proability or the equivalent event has measure zero.
\begin{definition}
    A set \(A \subset \Reals\) has Lebesgue measure zero if for every \(\epsilon > 0\), there exists a countable covering \(\set{A_i}\) of \(A\) by intervals such that 
    \begin{equation*}
        \sum_{i = 1}^{\infty} \func{\mu_L}{A_i} < \epsilon
    \end{equation*}
\end{definition}
\begin{itemize}
    \item Subset of a measure zero are measure zero.
    \item A signle point has a measure zero.
    \item countable union of measure zeros is a measure zero.
\end{itemize}

Let \(N = \set<\omega \in I>{ \frac{\func{s_N}{\omega}}{N} \to \frac{1}{2}}\). \(N\) is called the set of \textbf{normal numbers}. Let \(N^c\) be the complement of \(N\).
\begin{theorem}[Strong law of large numbers]
    SLLN states that \(N^c\) has measure zero.
\end{theorem}

\begin{proof}
    Let \(A_n = \set<\omega \in I>{\abs{\func{S_n}{\omega}} > \epsilon n}\) then 
    \begin{equation*}
        A_n = \set[[\Big]]<\omega \in I>{\bracket{\func{S_n}{\omega}}^4 > \epsilon^4 n^4}
    \end{equation*}
    By Chebyshev's inequality
    \begin{align*}
        \func{\mu_L}{A_n} &< \dfrac{1}{n^4 \epsilon^4} \int_0^1 \bracket{\func{S_n}{\omega}}^4 \diffOperator \omega \\
        &= \dfrac{1}{n^4 \epsilon^4} \bracket{n + 3n(n-1)} \leq \dfrac{3}{\epsilon^4 n^2}
    \end{align*}
    \begin{lemma}
        Given \(\delta > 0\) there exists a sequence \(\epsilon_1, \epsilon_2, \dots \) such that \(\epsilon \to 0\) and 
        \begin{equation*}
            \sum_{n = 1}^{\infty} \dfrac{3}{n^2 \epsilon_n^4} < \delta
        \end{equation*}
    \end{lemma}
    \begin{prooflemma}
        Choose \(\epsilon_n^4 = cn^{-1/2}\) for some constant \(c\). 
        \begin{equation*}
            \sum_{n = 1}^{\infty} \dfrac{3}{n^2 \epsilon_n^4} =\dfrac{3}{c} \sum_{n = 1}^{\infty} n^{-\frac{3}{2}} = \dfrac{3}{c}L < \delta
        \end{equation*}
        if \(c > \dfrac{3L}{\delta}\).
    \end{prooflemma}
    Finally, Let \(B_n = \set<\omega \in I> {\abs{\func{S_n}{\omega}} > \epsilon_n n}\). Then, by our first computation 
    \begin{equation*}
        \func{\mu_L}{B_n} < \dfrac{3}{\epsilon_n^4 n^2}
    \end{equation*}
    and by the lemma 
    \begin{equation*}
        \sum  \func{\mu_L}{B_n} < \dfrac{3}{\epsilon_n^4 n^2} < \delta 
    \end{equation*}
    It remains to show that \(N^c \subset \cup_{n = 1}^{\infty} B_n\) which is obvious. As every \(\omega \in N^c\) is some \(B_n\). Therefore, \(N^c\) has measure zero.
\end{proof}